\documentclass[journal,compsoc]{IEEEtran}

\usepackage[T1]{fontenc}
\usepackage[utf8]{inputenc}
\usepackage{microtype}
\usepackage{amsmath}
\usepackage{amssymb}
\usepackage{mathtools}
\usepackage{graphicx}
\usepackage{booktabs}
\usepackage{array}
\usepackage{float}
\usepackage{caption}
\usepackage{xcolor}
\usepackage{hyperref}
\usepackage{algorithm}
\usepackage{algpseudocode}
\usepackage{cite}

\hypersetup{
    colorlinks=true, linkcolor=blue, citecolor=blue, urlcolor=blue,
    pdftitle={Vel-v4.3: A Decentralized Swarm Protocol for Autonomous Volume Dominance},
    pdfauthor={Dakshith Rajkumar},
}

\newcommand{\vecc}[1]{\vec{\mathbf{#1}}}
\newcommand{\norm}[1]{\left\lVert#1\right\rVert}

\begin{document}

\title{Vel-v4.3: A Decentralized Swarm Protocol\\for Autonomous Volume Dominance}

\author{Dakshith~Rajkumar,~\IEEEmembership{Independent~Researcher}%
\IEEEcompsocitemizethanks{
\IEEEcompsocthanksitem D. Rajkumar is an Independent Researcher.
E-mail: daks2edu@gmail.com}}

\markboth{IEEE Transactions on Aerospace and Electronic Systems,~Vol.~X, No.~X, 2025}%
{Rajkumar: Vel-v4.3 Decentralized Swarm Protocol}

\maketitle

% ── Abstract ─────────────────────────────────────────────────────────────────
\begin{abstract}
This paper presents the comprehensive empirical validation of the Vel-v4.3 protocol,
a decentralized swarm coordination system for autonomous robotic navigation. A critical
architectural upgrade --- the replacement of center-outward bloom repulsion with
\textbf{Map-Aware Gradient Repulsion} (MAGR) --- is introduced and validated across
five systematic parameter sweeps comprising 105{,}000 simulation trials. The sweeps
characterize four independent protocol dimensions: agent population ($N$), repulsion
factor ($G$), momentum coefficient ($\mu$), and MAGR scan radius ($r$).

The Population Sweep ($N \in \{50,\,150,\,300\}$) confirms that MAGR restores monotonic
efficiency scaling eliminated by the original center-bloom architecture. The integer $G$
Sweep ($G \in \{0,\,2,\,4,\,6,\,8\}$) establishes $G=2$ as the coarse optimum and
confirms that repulsion accounts for over $97\%$ of protocol efficiency. The Fine $G$
Sweep ($G \in [1.0,\,3.0]$, $\Delta G=0.25$) refines the optimum to $G=1.5$ at $N=300$,
yielding $\bar{\eta}=1.796$. The Momentum Sweep ($\mu \in [0.70,\,0.99]$) identifies
$\mu=0.99$ as the empirical optimum, delivering $\bar{\eta}=2.696$ --- a $49.6\%$ gain
over the original $\mu=0.88$ assumption. The Scan Radius Sweep ($r \in [2,\,10]$) reveals
that MAGR efficiency is radius-invariant ($\Delta\eta=0.006$ across all tested values),
while computational cost scales by $17.6\times$ from $r=2$ to $r=10$, establishing $r=2$
as the optimal deployment radius.

The global efficiency optimum is $N=300$, $G=1.5$, $\mu=0.99$, $r=2$:
$\bar{\eta}=2.696$, $\sigma=0.089$, $\bar{T}=35$ steps --- a $78.7\%$ gain over the
initial baseline. The original operational benchmark of $\eta>0.105$ is superseded by
$\eta>2.0$ for fully-optimized MAGR configurations.
\end{abstract}

\begin{IEEEkeywords}
Swarm robotics, decentralized coordination, occupancy grid, map-aware repulsion,
spatio-temporal efficiency, momentum conservation, autonomous search, coverage path planning.
\end{IEEEkeywords}

% ── 1. Introduction ───────────────────────────────────────────────────────────
\section{Introduction}
\label{sec:intro}

\IEEEPARstart{I}{n} autonomous multi-agent robotics, the \textit{Redundant Search
Problem} describes the efficiency loss that arises when agents re-traverse previously
explored regions due to insufficient shared environmental awareness. Conventional swarm
architectures frequently rely on centralized command
structures~\cite{dorigo2004,brambilla2013,jamshidpey2024} --- an approach that introduces
a single point of failure, limits scalability, and prohibits genuine emergent behavior.

The Vel Protocol addresses both failure modes through a two-layer decentralized control
architecture in which global coverage efficiency emerges purely from local agent
interactions mediated by a shared spatial memory~\cite{reynolds1987}. Unlike centralized
systems, Vel requires no global coordinator: each agent acts independently, reading from
and writing to a shared binary occupancy grid that serves as the protocol's collective
intelligence layer.

This paper documents the most comprehensive empirical characterization of the Vel protocol
to date. Building on the introduction of Map-Aware Gradient Repulsion (MAGR), five
independent parameter sweeps are conducted across 105{,}000 simulation trials,
systematically characterizing efficiency across every major protocol dimension.

The central contributions are:
\begin{enumerate}
    \item Identification and correction of the boundary pile-up failure mode in
          center-outward bloom repulsion via MAGR.
    \item Empirical characterization of $\eta$ across 105{,}000 trials spanning
          four independent parameter dimensions.
    \item Discovery that the original momentum assumption ($\mu=0.88$) was
          suboptimal by $49.6\%$, with $\mu=0.99$ identified as the true optimum.
    \item Discovery that MAGR efficiency is radius-invariant, establishing $r=2$
          as the optimal deployment radius on computational cost grounds.
    \item Identification of the global efficiency optimum:
          $N=300$, $G=1.5$, $\mu=0.99$, $r=2$, $\bar{\eta}=2.696$.
    \item Establishment of a revised operational benchmark of $\eta>2.0$ for
          fully-optimized MAGR configurations.
\end{enumerate}

% ── 2. Related Work ───────────────────────────────────────────────────────────
\section{Related Work}
\label{sec:related}

\subsection{Coverage Path Planning}

Coverage Path Planning (CPP) is the task of determining a path that passes over all
points of an area or volume of interest while avoiding obstacles~\cite{galceran2013}.
A considerable body of research has addressed the CPP problem across domains including
vacuum cleaning robots, autonomous underwater vehicles, demining robots, and aerial
survey platforms. Classical CPP approaches such as boustrophedon decomposition partition
the environment into cells and plan optimal traversal sequences, but are designed for
single agents and do not address the coordination challenges that arise when multiple
agents cover the same space.

\subsection{Ant Colony Optimization}

Ant Colony Optimization (ACO), originally introduced by Dorigo~\cite{dorigo1992} for
the Traveling Salesman Problem, models the pheromone-trail communication of biological
ants to solve combinatorial optimization problems. ACO has been extensively applied to
multi-robot path planning, where pheromone trails guide agents toward unexplored regions.
The key limitation of ACO for coverage tasks is stagnation: once pheromone trails
converge on a particular traversal pattern, agents are difficult to redirect toward
remaining unexplored regions. Recent work has addressed this through adaptive pheromone
update rules~\cite{li2024}, but the fundamental reliance on pheromone convergence
constrains real-time adaptability.

Vel differs from ACO in a fundamental architectural choice: rather than pheromone
concentration, Vel uses a binary occupancy grid in which cells are either explored or
unexplored with no intermediate state. This eliminates the stagnation problem entirely
--- a cell is explored or it is not --- and removes the need for pheromone decay
parameter tuning.

\subsection{Particle Swarm Optimization}

Particle Swarm Optimization (PSO), introduced by Kennedy and
Eberhart~\cite{kennedy1995}, models the collective motion of flocking agents to
optimize continuous objective functions. PSO has been adapted for multi-robot coverage
tasks, where agents are attracted toward unexplored regions by a shared global best
position. However, PSO-based coverage systems are vulnerable to premature convergence:
once the global best position is identified, all agents converge toward it rather than
distributing across the remaining unexplored space.

The Vel KVL subsystem draws conceptual inspiration from PSO's velocity update rule ---
specifically the momentum term $\mu$ that damps high-frequency direction changes. However,
Vel replaces PSO's global attractor with a map-aware local gradient, avoiding the
premature convergence failure mode by ensuring each agent steers toward the nearest
\emph{locally} unexplored cell rather than a globally optimal position.

\subsection{Decentralized Multi-Robot Coverage}

The problem of centralization versus decentralization in multi-robot sweep coverage has
been directly studied by Jamshidpey et al.~\cite{jamshidpey2024}, who demonstrate that
decentralized approaches achieve competitive coverage efficiency while eliminating the
single point of failure inherent in centralized command architectures. Brambilla et
al.~\cite{brambilla2013} provide a comprehensive review of the swarm engineering
perspective, cataloging the conditions under which emergent decentralized behavior
outperforms designed centralized coordination. The Vel protocol is positioned squarely
within this decentralized paradigm, extending it with a formally defined shared memory
layer that enables implicit coordination without explicit communication.

\subsection{Occupancy Grid Representations}

The occupancy grid representation used by Vel's SSM layer was introduced by Moravec and
Elfes~\cite{moravec1985} and formalized by Elfes~\cite{elfes1989} for mobile robot
perception. The original formulation used probabilistic cell states to represent sensor
uncertainty. Vel simplifies this to a deterministic binary representation (explored/unexplored)
appropriate for the coverage task, where the question of interest is not the probability
of occupancy but the certainty of prior visitation.

\subsection{Positioning of Vel}

The Vel protocol occupies a distinct position in the swarm coverage landscape. Unlike
ACO, it uses deterministic binary memory rather than probabilistic pheromones. Unlike
PSO, it uses local map-aware gradient descent rather than global attraction. Unlike
classical CPP methods, it is fully decentralized with no pre-planned traversal sequence.
The MAGR architecture, introduced in this paper, is a novel contribution that addresses
the boundary pile-up failure mode not previously identified or resolved in the
decentralized coverage literature.

% ── 3. System Architecture ────────────────────────────────────────────────────
\section{System Architecture}
\label{sec:arch}

The Vel protocol operates as a two-layer control system. The lower layer ---
the \textbf{Shared Spatial Memory} (SSM) --- maintains a globally synchronized binary
map of the environment that every agent reads and writes simultaneously, functioning as
a distributed hive-mind. The upper layer --- the \textbf{Kinetic Vector Logic} (KVL)
--- governs individual agent motion through continuous force integration at 50\,Hz,
using the SSM to continuously steer each agent toward undiscovered territory.

\subsection{Shared Spatial Memory (SSM)}

The operational theater is discretized into an occupancy grid $M$ of dimensions
$W \times H$~\cite{moravec1985,elfes1989}. Each cell $c \in M$ is a binary state
variable:
\begin{equation}
    c(x,y) = \begin{cases} 0 & \text{if Unexplored} \\ 1 & \text{if Explored}
    \end{cases}
    \label{eq:cell}
\end{equation}
Agents perform atomic \texttt{write(1)} operations on every visited cell. The SSM
acts as a globally shared repulsion field: agents query it continuously to identify
and prioritize \emph{White Space} (cells where $c=0$). Because every agent reads the
same SSM, the system achieves implicit coordination without any direct agent-to-agent
communication channel --- a key architectural advantage over ACO and PSO variants that
require explicit message passing.

Coverage is defined as the fraction of cells with $c=1$:
\begin{equation}
    \text{Coverage}(t) = \frac{\sum_{x,y} c(x,y,t)}{W \cdot H}
\end{equation}
Each trial terminates when $\text{Coverage}(t) \geq 0.95$ or when the step budget
$T_{\max}=1{,}000$ is exhausted.

\subsection{Kinetic Vector Logic (KVL)}

Agent $i$ maintains a velocity vector $\vecc{V}_i = [V_x,\,V_y]^\top$, updated at
50\,Hz:
\begin{equation}
    \vecc{V}_{i}(t+1) = \Bigl(\vecc{V}_{i}(t) + \vecc{R}_G(\theta)\Bigr)\cdot\mu
    \label{eq:kvl}
\end{equation}
where $\vecc{R}_G(\theta) = G\cdot[\cos\theta,\,\sin\theta]^\top\cdot 0.12$ is the
Repulsion Vector derived from the SSM gradient, $\mu \in (0,\,1]$ is the Momentum
Conservation Coefficient governing velocity decay, and $G \in \mathbb{R}^+$ is the
tunable Repulsion Factor.

The momentum term $\mu$ serves an analogous role to the inertia weight in PSO~\cite{kennedy1995}:
it prevents high-frequency direction oscillations and maintains smooth, arc-shaped
agent trajectories. The Momentum Sweep (Study IV) provides the first systematic
empirical characterization of $\mu$'s effect on coverage efficiency.

\subsection{Map-Aware Gradient Repulsion (MAGR)}

\subsubsection{Failure Mode of Center-Outward Bloom}
The original Vel implementation computed $\theta$ as the bearing from the swarm's
geometric center to agent $i$'s position. This produced a systematic failure for
square grids: center-bloom drives agents radially outward in a circle, but the four
corners of a square lie significantly further from center than edge midpoints. Agents
reaching the perimeter are repelled back toward already-explored space rather than
redirected toward nearby unexplored corners --- the \textit{boundary pile-up effect}.
Empirically, this caused $N=300$ to underperform $N=150$, as additional agents
compounded the pile-up rather than contributing new coverage. This failure mode is
analogous to the stagnation problem in ACO, where positive feedback drives agents
toward already-converged paths.

\subsubsection{MAGR Algorithm}
MAGR replaces the center-bearing calculation with a local SSM neighborhood scan.
At each timestep, agent $i$ scans a radius-$r$ Moore neighborhood and steers toward
the nearest unexplored cell:
\begin{equation}
    \vecc{R}_G = \begin{cases}
        G \cdot \dfrac{(dx^*,\,dy^*)}{\norm{(dx^*,\,dy^*)}}\cdot 0.12
            & \text{unexplored cell found within radius }r \\[0.8em]
        G \cdot \hat{u}_{\text{out}}\cdot 0.12
            & \text{fallback: center-bloom}
    \end{cases}
    \label{eq:magr}
\end{equation}
where:
\begin{equation}
    (dx^*,\,dy^*) = \underset{\substack{(dx,dy):\,\norm{(dx,dy)}\leq r \\ c(g_x+dx,\,g_y+dy)=0}}{\arg\min}\;\norm{(dx,dy)}
\end{equation}

\begin{algorithm}[H]
\caption{MAGR Repulsion Vector Computation for Agent $i$}
\label{alg:magr}
\begin{algorithmic}[1]
\State $(g_x,g_y)\gets\lfloor(p_x^i,p_y^i)\rfloor$
\If{$c(g_x,g_y)=0$} $c(g_x,g_y)\gets 1$ \EndIf
\State $d^*\gets\infty$; \texttt{found}$\gets$\textbf{false}
\For{$dy\in[-r,r]$} \textbf{ do}
    \For{$dx\in[-r,r]$} \textbf{ do}
        \If{$(g_x{+}dx,\,g_y{+}dy)\in M$ \textbf{and}
            $c(g_x{+}dx,\,g_y{+}dy)=0$ \textbf{and}
            $\sqrt{dx^2{+}dy^2}<d^*$}
            \State $(dx^*,dy^*)\gets(dx,dy)$; $d^*\gets\sqrt{dx^2{+}dy^2}$;
                   \texttt{found}$\gets$\textbf{true}
        \EndIf
    \EndFor
\EndFor
\If{\texttt{found}}
    \State \Return $G\cdot(dx^*,dy^*)/\norm{(dx^*,dy^*)}\cdot 0.12$
\Else
    \State \Return $G\cdot[\cos\theta_{\text{out}},\sin\theta_{\text{out}}]^\top\cdot 0.12$
\EndIf
\end{algorithmic}
\end{algorithm}

% ── 4. Performance Metrics ────────────────────────────────────────────────────
\section{Performance Metrics}
\label{sec:metrics}

\subsection{Spatio-Temporal Efficiency ($\eta$)}

The primary performance metric is Spatio-Temporal Efficiency:
\begin{equation}
    \eta = \frac{\Delta C}{\Delta T \cdot N}
    \label{eq:eta}
\end{equation}
where $\Delta C$ is newly explored cells, $\Delta T$ elapsed timesteps, $N$ agent count.
This metric captures the joint optimization objective: maximize coverage, minimize time,
minimize agents. A higher $\eta$ indicates more efficient use of both time and agents.
Note that $\eta$ is normalized by $N$, so a larger swarm must achieve proportionally
more coverage to maintain the same $\eta$ --- this normalization reveals diminishing
returns and scan-overlap congestion that raw coverage rates would conceal.

\subsection{Stability ($\sigma$)}

The trial-to-trial standard deviation of $\eta$:
\begin{equation}
    \sigma = \sqrt{\frac{1}{K}\sum_{k=1}^{K}(\eta_k - \bar{\eta})^2}
\end{equation}
where $K$ is the number of trials. Lower $\sigma$ indicates more deterministic,
hardware-deployable behavior. The efficiency--stability tradeoff identified in Studies
II through IV has direct implications for deployment: configurations with high
$\bar{\eta}$ but high $\sigma$ may be unsuitable for safety-critical applications
requiring predictable coverage completion times.

\subsection{Coordination Quality ($\eta/N$)}

Per-agent efficiency, revealing diminishing returns as $N$ scales:
\begin{equation}
    Q = \frac{\eta}{N}
\end{equation}
Decreasing $Q$ with increasing $N$ indicates the onset of scan-overlap congestion,
where multiple agents simultaneously target the same unexplored cells.

\subsection{Convergence Speed ($\bar{T}$)}

Mean timesteps to $95\%$ grid coverage across all trials. For $G=0$ configurations
where coverage is never achieved, $\bar{T}$ is reported as the step budget ceiling
($T_{\max}=1{,}000$).

% ── 5. Study I: Population Sweep ─────────────────────────────────────────────
\section{Study I: Population Sweep}
\label{sec:nsweep}

\subsection{Configuration}
Three agent populations at fixed $G=4.0$, $\mu=0.88$, $r=5$. 5{,}000 trials each on
a $50\times50$ grid (15{,}000 total).

\begin{table}[H]
\renewcommand{\arraystretch}{1.3}
\caption{Population Sweep Results (5,000 trials, $G=4.0$, $\mu=0.88$, $r=5$)}
\label{tab:nsweep}
\centering
\begin{tabular}{ccccc}
\hline
$N$ & $\bar{\eta}$ & $\sigma_\eta$ & $\bar{T}$ & $Q\times10^3$ \\
\hline
50  & 0.4345 & 0.04026 & 247 & 8.69 \\
150 & 0.9076 & 0.07062 & 106 & 6.05 \\
300 & 1.5092 & 0.11091 & 63  & 5.03 \\
\hline
\end{tabular}
\end{table}

\begin{figure}[H]
\centering
\includegraphics[width=\columnwidth]{N_Sweep_Data_Visualization/fig1_convergence.png}
\caption{Rolling mean $\eta$ ($w=100$) across 5{,}000 trials per population. All
         configurations stabilize rapidly. Shaded bands: 95\% CI.}
\label{fig:n_convergence}
\end{figure}

\begin{figure}[H]
\centering
\includegraphics[width=\columnwidth]{N_Sweep_Data_Visualization/fig5_summary_bars.png}
\caption{Population Sweep summary. Clockwise from top-left: mean $\eta\pm1\sigma$,
         stability $\sigma$, Coordination Quality, mean convergence steps.}
\label{fig:n_summary}
\end{figure}

\begin{figure}[H]
\centering
\includegraphics[width=\columnwidth]{N_Sweep_Data_Visualization/fig6_scaling_laws.png}
\caption{Left: power law fit $\bar{\eta}\propto N^\alpha$ with 95\% CI ribbon.
         Right: convergence steps vs.\ $N$ with linear fit.}
\label{fig:n_scaling}
\end{figure}

\begin{figure}[H]
\centering
\includegraphics[width=0.85\columnwidth]{N_Sweep_Data_Visualization/fig10_per_agent.png}
\caption{Coordination Quality $Q$ vs.\ $N$. Monotonic decrease confirms scan-overlap
         congestion onset. Error bars: $\pm1\sigma$.}
\label{fig:n_per_agent}
\end{figure}

\subsection{Analysis}

\paragraph{Monotonic Scaling Under MAGR.}
MAGR produces strictly monotonic $N$-scaling: $+108.9\%$ ($N$: $50\to150$) and
$+66.3\%$ ($N$: $150\to300$). This is the primary validation of the MAGR upgrade ---
the original center-bloom implementation exhibited a congestion inversion where $N=300$
underperformed $N=150$ due to boundary pile-up. MAGR eliminates this inversion entirely.

\paragraph{Power Law Scaling.}
The relationship $\bar{\eta}\propto N^\alpha$ holds across the tested range, enabling
performance prediction at deployment densities beyond those tested. This power law is
consistent with theoretical predictions for decentralized coverage systems where agent
interactions are mediated by a shared memory field~\cite{brambilla2013}.

\paragraph{Diminishing Returns Per Agent.}
Coordination Quality $Q$ decreases monotonically: $8.69\times10^{-3}$ ($N=50$),
$6.05\times10^{-3}$ ($N=150$), $5.03\times10^{-3}$ ($N=300$). Each additional agent
delivers less marginal coverage gain, consistent with scan-overlap congestion where
multiple agents simultaneously target the same unexplored cell. The theoretical optimal
density ceiling --- where adding agents yields negative marginal returns on $\eta$ ---
is not reached within the tested range, suggesting further gains are achievable at
higher $N$.

% ── 6. Study II: Integer G Sweep ─────────────────────────────────────────────
\section{Study II: Integer Repulsion Factor Sweep}
\label{sec:gsweep}

\subsection{Configuration}
Five $G$ values at $N\in\{150,300\}$, including $G=0$ control. 2{,}000 trials each
(20{,}000 total). $\mu=0.88$, $r=5$ fixed.

\begin{table}[H]
\renewcommand{\arraystretch}{1.3}
\caption{Integer $G$ Sweep Results ($^\dagger$budget exhausted; $\Delta\eta$ vs.\ $G{=}4$)}
\label{tab:gsweep}
\centering
\begin{tabular}{lcccccc}
\hline
Config & $N$ & $G$ & $\bar{\eta}$ & $\sigma_\eta$ & $\bar{T}$ & $\Delta\eta$ \\
\hline
G0-N150 & 150 & 0 & 0.0262 & 0.00057 & $1000^\dagger$ & $-97.1\%$ \\
G2-N150 & 150 & 2 & 0.9161 & 0.07931 & 105 & $+0.9\%$ \\
G4-N150 & 150 & 4 & 0.9076 & 0.07135 & 106 & baseline \\
G6-N150 & 150 & 6 & 0.8727 & 0.06666 & 110 & $-3.8\%$ \\
G8-N150 & 150 & 8 & 0.8821 & 0.06805 & 109 & $-2.8\%$ \\
\hline
G0-N300 & 300 & 0 & 0.0284 & 0.00020 & $1000^\dagger$ & $-98.1\%$ \\
G2-N300 & 300 & 2 & 1.7732 & 0.17188 & 54  & $\mathbf{+17.5\%}$ \\
G4-N300 & 300 & 4 & 1.5087 & 0.11017 & 64  & baseline \\
G6-N300 & 300 & 6 & 1.4261 & 0.09883 & 67  & $-5.5\%$ \\
G8-N300 & 300 & 8 & 1.4909 & 0.08969 & 64  & $-1.2\%$ \\
\hline
\end{tabular}
\end{table}

\begin{figure}[H]
\centering
\includegraphics[width=\columnwidth]{G_Sweep_Data_Visualization/figG9_control_G0.png}
\caption{$G=0$ vs.\ best $G$ per population. Without repulsion, $\eta$ collapses by
         $97$--$98\%$, confirming repulsion is the primary protocol efficiency driver.}
\label{fig:g0_control}
\end{figure}

\begin{figure}[H]
\centering
\includegraphics[width=\columnwidth]{G_Sweep_Data_Visualization/figG2_eta_vs_G.png}
\caption{Mean $\eta\pm1\sigma$ vs.\ integer $G$ for $N=150$ and $N=300$.
         $G=2$ is the coarse optimum for both populations.}
\label{fig:g_eta}
\end{figure}

\begin{figure}[H]
\centering
\includegraphics[width=\columnwidth]{G_Sweep_Data_Visualization/figG6_heatmap.png}
\caption{$N\times G$ efficiency heatmap. Left: mean $\eta$. Right: mean steps.
         Global coarse optimum: $N=300$, $G=2$.}
\label{fig:g_heatmap}
\end{figure}

\begin{figure}[H]
\centering
\includegraphics[width=0.9\columnwidth]{G_Sweep_Data_Visualization/figG10_ranking.png}
\caption{All integer $G$ configurations ranked by $\bar{\eta}$. $N=300$ dominates
         the top four positions regardless of $G$.}
\label{fig:g_ranking}
\end{figure}

\subsection{Analysis}

\paragraph{Repulsion is Non-Negotiable.}
The $G=0$ control establishes a critical baseline: without repulsion, Vel degenerates
to an unguided random walk, achieving $\eta \approx 0.027$ --- a $97$--$98\%$
efficiency collapse relative to the $G=4$ baseline. This result empirically confirms
the repulsion subsystem as the primary driver of protocol efficiency and validates
the MAGR architecture as essential rather than optional.

\paragraph{Overshoot Mechanism.}
$G=4$ underperforms $G=2$ because excessive repulsion force causes agents to overshoot
target unexplored cells. When an agent's MAGR vector points toward a nearby unexplored
cell at distance $d$, the repulsion impulse $G\cdot 0.12$ must be proportioned to
that distance. At $G=4$, the impulse frequently exceeds what is required to reach the
target, carrying the agent past it and requiring corrective steps. At $G=2$, the
impulse is calibrated to the typical MAGR target distance, minimizing overshoot.

\paragraph{Stability--Efficiency Tradeoff.}
$\sigma$ decreases monotonically with $G$ at $N=300$: from $0.172$ ($G=2$) to $0.090$
($G=8$) --- a $47.7\%$ variance reduction at the cost of $15.9\%$ mean efficiency.
Hardware deployments requiring deterministic coverage completion times may prefer
higher $G$ values despite the efficiency penalty. This tradeoff is formally quantified
in Study IV.

\paragraph{Integer Resolution Insufficient.}
The coarse $G$ sweep identifies $G=2$ as the optimum but tests only integer values.
The true efficiency peak may lie between $G=1$ and $G=2$, motivating the fine sweep
in Study III.

% ── 7. Study III: Fine G Sweep ───────────────────────────────────────────────
\section{Study III: Fine Repulsion Factor Sweep}
\label{sec:fgsweep}

\subsection{Configuration}
Nine $G$ values in $[1.0,\,3.0]$ at $\Delta G=0.25$, for $N\in\{150,\,300\}$.
2{,}000 trials each (36{,}000 total). $\mu=0.88$, $r=5$ fixed.

\begin{table}[H]
\renewcommand{\arraystretch}{1.3}
\caption{Fine $G$ Sweep Results (2,000 trials each, $\mu=0.88$, $r=5$)}
\label{tab:fgsweep}
\centering
\begin{tabular}{ccccccc}
\hline
$G$ & \multicolumn{3}{c}{$N=150$} & \multicolumn{3}{c}{$N=300$} \\
    & $\bar{\eta}$ & $\sigma$ & $\bar{T}$ & $\bar{\eta}$ & $\sigma$ & $\bar{T}$ \\
\hline
1.00 & 0.8970 & 0.108 & 108 & 1.6394 & 0.122 & 58 \\
1.25 & 0.9130 & 0.114 & 106 & 1.7415 & 0.135 & 55 \\
1.50 & 0.9097 & 0.099 & 106 & \textbf{1.7958} & 0.159 & \textbf{53} \\
1.75 & 0.9123 & 0.085 & 105 & 1.7788 & 0.167 & 54 \\
2.00 & 0.9174 & 0.080 & 105 & 1.7714 & 0.175 & 54 \\
2.25 & 0.9137 & 0.077 & 105 & 1.7069 & 0.179 & 56 \\
2.50 & 0.9208 & 0.074 & 104 & 1.6317 & 0.170 & 59 \\
2.75 & \textbf{0.9222} & 0.076 & 104 & 1.6022 & 0.152 & 60 \\
3.00 & 0.9153 & 0.072 & 105 & 1.5373 & 0.128 & 62 \\
\hline
\end{tabular}
\end{table}

\begin{figure}[H]
\centering
\includegraphics[width=\columnwidth]{Fine_G_Sweep/figFG1_eta_vs_G.png}
\caption{Fine $G$ sweep: mean $\eta\pm1\sigma$ vs.\ $G$. Peaks: $G=2.75$ ($N=150$),
         $G=1.5$ ($N=300$). Dashed lines mark empirical optima per population.}
\label{fig:fg_eta}
\end{figure}

\begin{figure}[H]
\centering
\includegraphics[width=\columnwidth]{Fine_G_Sweep/figFG4_comparison.png}
\caption{$N=150$ vs.\ $N=300$ efficiency across fine $G$ sweep. Peak locations
         diverge, confirming a $N\times G$ interaction effect.}
\label{fig:fg_comparison}
\end{figure}

\begin{figure}[H]
\centering
\includegraphics[width=\columnwidth]{Fine_G_Sweep/figFG5_heatmap.png}
\caption{Fine $G$ sweep $N\times G$ heatmap. Left: mean $\eta$. Right: mean steps.
         Color gradient reveals the interaction structure.}
\label{fig:fg_heatmap}
\end{figure}

\begin{figure}[H]
\centering
\includegraphics[width=\columnwidth]{Fine_G_Sweep/figFG8_ranking.png}
\caption{All fine $G$ configurations ranked by $\bar{\eta}$. $N=300$, $G=1.5$
         is the fine-sweep global optimum.}
\label{fig:fg_ranking}
\end{figure}

\subsection{Analysis}

\paragraph{Population-Dependent Optimal G.}
The fine sweep reveals that the optimal repulsion factor is population-dependent:
$G^*_{N=300}=1.5$ versus $G^*_{N=150}=2.75$. This constitutes a genuine $N\times G$
interaction effect: higher populations require lower repulsion to avoid the overshoot
mechanism identified in Study II, because the mean inter-agent distance is smaller
at higher density, and the MAGR target cells are closer on average. Lower repulsion
force is therefore sufficient to reach them cleanly. At $N=150$, agents are more
dispersed, and a stronger repulsion impulse is needed to maintain outward pressure.

\paragraph{Revised G Optimum.}
The true optimum for $N=300$ is $G=1.5$ ($\bar{\eta}=1.796$), not $G=2$ as the
integer sweep suggested. The difference is $1.4\%$ in mean efficiency --- modest but
consistent across 2{,}000 trials and therefore statistically reliable.

\paragraph{Stability Gradient.}
$\sigma$ decreases monotonically with $G$ for $N=300$: higher repulsion produces more
deterministic behavior. This creates a practical design space: engineers can select
any $G \in [1.5,\,3.0]$ and trade efficiency for stability in a controlled manner.

% ── 8. Study IV: Momentum Sweep ──────────────────────────────────────────────
\section{Study IV: Momentum Coefficient Sweep}
\label{sec:musweep}

\subsection{Configuration}
Eight $\mu$ values in $[0.70,\,0.99]$ at $N=300$, $G=1.5$, $r=5$.
2{,}000 trials each (16{,}000 total).

\begin{table}[H]
\renewcommand{\arraystretch}{1.3}
\caption{Momentum Sweep Results ($N=300$, $G=1.5$, $r=5$)}
\label{tab:musweep}
\centering
\begin{tabular}{ccccc}
\hline
$\mu$ & $\bar{\eta}$ & $\sigma_\eta$ & $\bar{T}$ & $\Delta\eta$ vs.\ $\mu{=}0.88$ \\
\hline
0.70 & 0.7574 & 0.05768 & 126 & $-58.0\%$ \\
0.75 & 0.9650 & 0.06852 &  99 & $-46.5\%$ \\
0.80 & 1.2226 & 0.09524 &  78 & $-32.2\%$ \\
0.85 & 1.5481 & 0.12655 &  62 & $-14.1\%$ \\
0.88 & 1.8026 & 0.16393 &  53 & baseline \\
0.90 & 1.9071 & 0.17778 &  50 & $+5.8\%$ \\
0.95 & 2.3871 & 0.18746 &  40 & $+32.4\%$ \\
\textbf{0.99} & \textbf{2.6963} & \textbf{0.08897} & \textbf{35} & $\mathbf{+49.6\%}$ \\
\hline
\end{tabular}
\end{table}

\begin{figure}[H]
\centering
\includegraphics[width=\columnwidth]{Momentum_Sweep/figMU1_eta_vs_mu.png}
\caption{Mean $\eta\pm1\sigma$ (left) and convergence speed (right) vs.\ $\mu$.
         Both metrics improve monotonically across $[0.70,\,0.99]$.}
\label{fig:mu_eta}
\end{figure}

\begin{figure}[H]
\centering
\includegraphics[width=\columnwidth]{Momentum_Sweep/figMU6_gain.png}
\caption{\% efficiency gain over original $\mu=0.88$ baseline. Red bars indicate
         configurations worse than the original; green bars indicate improvements.}
\label{fig:mu_gain}
\end{figure}

\begin{figure}[H]
\centering
\includegraphics[width=\columnwidth]{Momentum_Sweep/figMU7_boxplots.png}
\caption{Full $\eta$ distribution shapes across $\mu$ values. Note the sharp
         variance reduction at $\mu=0.99$ relative to $\mu=0.95$.}
\label{fig:mu_boxplots}
\end{figure}

\begin{figure}[H]
\centering
\includegraphics[width=\columnwidth]{Momentum_Sweep/figMU8_frontier.png}
\caption{Efficiency--stability frontier. $\mu=0.99$ (top-left) occupies the dominant
         position: simultaneously highest $\bar{\eta}$ and lowest $\sigma$.}
\label{fig:mu_frontier}
\end{figure}

\subsection{Analysis}

\paragraph{The Original $\mu=0.88$ Was Suboptimal.}
The momentum sweep is the most consequential result of this study. The original
$\mu=0.88$ was an untested assumption carried forward from the initial protocol
specification. The sweep reveals it was suboptimal by $49.6\%$: $\mu=0.99$ achieves
$\bar{\eta}=2.696$ versus $\bar{\eta}=1.803$ at $\mu=0.88$.

\paragraph{Monotonic Efficiency Scaling.}
Unlike the $G$ sweep, the momentum sweep reveals no interior optimum within the tested
range: efficiency increases monotonically from $\mu=0.70$ to $\mu=0.99$. This suggests
the true optimum may lie at or near $\mu=1.0$, i.e., zero momentum decay. While
$\mu=1.0$ is physically unrealizable on real hardware (all physical systems exhibit
some momentum loss), $\mu=0.99$ is practically achievable and provides $49.6\%$
efficiency gains over the original assumption.

\paragraph{Dominant Frontier Position.}
$\mu=0.99$ simultaneously achieves the highest mean efficiency ($\bar{\eta}=2.696$)
and the lowest variance ($\sigma=0.089$) of any tested configuration. This is a rare
result in parameter optimization: typically, maximum efficiency comes at the cost of
increased variance. The mechanism is that high momentum causes agents to maintain
smooth, consistent sweeping arcs rather than erratic directional changes, reducing
both wasted steps and trial-to-trial variability.

\paragraph{Variance Spike at $\mu=0.95$.}
A notable non-monotonicity appears in the stability data: $\sigma$ increases from
$0.164$ ($\mu=0.88$) to $0.187$ ($\mu=0.95$) before dropping sharply to $0.089$
($\mu=0.99$). This suggests a phase transition in agent trajectory dynamics between
$\mu=0.95$ and $\mu=0.99$, where the system transitions from a high-variance
exploratory regime to a low-variance stable sweeping regime.

% ── 9. Study V: Scan Radius Sweep ────────────────────────────────────────────
\section{Study V: Scan Radius Sweep}
\label{sec:rsweep}

\subsection{Configuration}
Nine $r$ values in $[2,\,10]$ at $N=300$, $G=1.5$, $\mu=0.99$.
2{,}000 trials each (18{,}000 total).

\begin{table}[H]
\renewcommand{\arraystretch}{1.3}
\caption{Scan Radius Sweep Results ($N=300$, $G=1.5$, $\mu=0.99$)}
\label{tab:rsweep}
\centering
\begin{tabular}{ccccc}
\hline
$r$ & $\bar{\eta}$ & $\sigma_\eta$ & $\bar{T}$ & Reads/agent/step \\
\hline
2  & 2.6892 & 0.09360 & 35.5 & 25 \\
3  & 2.6924 & 0.09102 & 35.4 & 49 \\
4  & 2.6948 & 0.09536 & 35.4 & 81 \\
5  & 2.6909 & 0.09255 & 35.4 & 121 \\
6  & 2.6916 & 0.09301 & 35.4 & 169 \\
7  & 2.6943 & 0.09297 & 35.4 & 225 \\
8  & 2.6956 & 0.09163 & 35.4 & 289 \\
9  & 2.6955 & 0.09245 & 35.4 & 361 \\
10 & 2.6922 & 0.09373 & 35.4 & 441 \\
\hline
\end{tabular}
\end{table}

\begin{figure}[H]
\centering
\includegraphics[width=\columnwidth]{Radius_Sweep/figR1_eta_vs_r.png}
\caption{Mean $\eta\pm1\sigma$ (left) and convergence steps (right) vs.\ $r$.
         Efficiency is flat across all tested radii --- the radius-invariance result.}
\label{fig:r_eta}
\end{figure}

\begin{figure}[H]
\centering
\includegraphics[width=\columnwidth]{Radius_Sweep/figR7_cost_vs_eta.png}
\caption{Computational cost (SSM reads/agent/step) vs.\ $r$, overlaid with $\bar{\eta}$.
         Cost scales as $(2r+1)^2$ while efficiency remains flat: a $17.6\times$ cost
         increase for $\Delta\eta=0.006$.}
\label{fig:r_cost}
\end{figure}

\begin{figure}[H]
\centering
\includegraphics[width=\columnwidth]{Radius_Sweep/figR6_boxplots.png}
\caption{Full $\eta$ distribution shapes across all $r$ values. Near-identical
         distributions confirm radius-invariance.}
\label{fig:r_boxplots}
\end{figure}

\begin{figure}[H]
\centering
\includegraphics[width=\columnwidth]{Radius_Sweep/figR8_heatmap.png}
\caption{Scan radius sweep summary heatmap: $\bar{\eta}$, $\sigma$, and $\bar{T}$
         across all $r$. Near-uniform color confirms invariance.}
\label{fig:r_heatmap}
\end{figure}

\subsection{Analysis}

\paragraph{MAGR is Radius-Invariant.}
The scan radius sweep produces the most operationally significant finding of the study.
MAGR efficiency is radius-invariant: $\Delta\eta=0.006$ across $r\in[2,\,10]$, which
is within the normal trial-to-trial variance of any single configuration. Convergence
speed is similarly invariant: $\bar{T}$ differs by at most 0.1 steps across all tested
radii.

\paragraph{Computational Cost Implication.}
MAGR scan cost scales as $(2r+1)^2$: from 25 reads per agent per step at $r=2$ to 441
reads at $r=10$, a $17.6\times$ increase for zero efficiency gain. The optimal
deployment radius is therefore $r=2$: minimum scan bandwidth, maximum deployability,
identical performance to any larger radius.

\paragraph{Theoretical Explanation.}
The radius-invariance result provides theoretical insight into MAGR's mechanism. At
the optimal configuration ($\mu=0.99$, $G=1.5$), agents maintain smooth sweeping arcs
that carry them continuously toward frontier boundaries. The nearest unexplored cell is
almost always within 2 cells of the agent's current position during active coverage,
making larger scan radii redundant. The MAGR algorithm effectively degenerates to a
radius-1 or radius-2 scan at runtime regardless of the configured radius, because
there is almost always an unexplored cell within that distance during the coverage
process. Only near the very end of a trial, when coverage approaches $95\%$, would
a larger scan radius provide any practical advantage --- and the efficiency metric
integrates over the entire trial, diluting this effect.

% ── 10. Combined Results ─────────────────────────────────────────────────────
\section{Combined Results and Global Optimum}
\label{sec:combined}

\subsection{Evolution of the Optimum}

Each successive study refined the protocol optimum. Table~\ref{tab:evolution} traces
this evolution from the initial baseline through all five studies.

\begin{table}[H]
\renewcommand{\arraystretch}{1.3}
\caption{Evolution of the Protocol Optimum Across Studies}
\label{tab:evolution}
\centering
\begin{tabular}{llccc}
\hline
\textbf{After Study} & \textbf{Best Config} & $\bar{\eta}$ & $\sigma$ & $\bar{T}$ \\
\hline
I  (N sweep)      & N=300, G=4.0, $\mu$=0.88, r=5 & 1.509 & 0.111 & 63 \\
II (int. G sweep) & N=300, G=2.0, $\mu$=0.88, r=5 & 1.773 & 0.172 & 54 \\
III (fine G)      & N=300, G=1.5, $\mu$=0.88, r=5 & 1.796 & 0.159 & 53 \\
IV ($\mu$ sweep)      & N=300, G=1.5, $\mu$=0.99, r=5 & 2.696 & 0.089 & 35 \\
V  (r sweep)      & N=300, G=1.5, $\mu$=0.99, r=2 & 2.689 & 0.094 & 36 \\
\hline
\end{tabular}
\end{table}

The Study IV momentum optimization dominates all other gains: the $\mu$ upgrade alone
accounts for $+49.6\%$ efficiency improvement, compared to $+17.4\%$ across all $G$
and $N$ optimization. This establishes $\mu$ as the single most impactful tunable
parameter in the Vel-v4.3 protocol.

\subsection{Master Configuration Ranking}

Table~\ref{tab:master} presents the complete ranking of all configurations tested
across Studies II through V, ordered by $\bar{\eta}$.

\begin{table}[H]
\renewcommand{\arraystretch}{1.2}
\caption{Master Configuration Ranking (All Studies)}
\label{tab:master}
\centering
\begin{tabular}{clcccc}
\hline
Rank & Config & $\bar{\eta}$ & $\sigma$ & $\bar{T}$ & Study \\
\hline
1  & N300, G1.5, $\mu$0.99, r=5 & 2.696 & 0.089 & 35 & IV \\
2  & N300, G1.5, $\mu$0.99, r=8 & 2.696 & 0.092 & 35 & V \\
3  & N300, G1.5, $\mu$0.99, r=9 & 2.696 & 0.092 & 35 & V \\
4  & N300, G1.5, $\mu$0.99, r=4 & 2.695 & 0.095 & 35 & V \\
5  & N300, G1.5, $\mu$0.99, r=7 & 2.694 & 0.093 & 35 & V \\
6  & N300, G1.5, $\mu$0.99, r=6 & 2.692 & 0.093 & 35 & V \\
7  & N300, G1.5, $\mu$0.99, r=3 & 2.692 & 0.091 & 35 & V \\
8  & N300, G1.5, $\mu$0.99, r=10& 2.692 & 0.094 & 35 & V \\
9  & N300, G1.5, $\mu$0.99, r=2 & 2.689 & 0.094 & 36 & V \\
10 & N300, G1.5, $\mu$0.95, r=5 & 2.387 & 0.187 & 40 & IV \\
11 & N300, G1.5, $\mu$0.90, r=5 & 1.907 & 0.178 & 50 & IV \\
12 & N300, G1.5, $\mu$0.88, r=5 & 1.803 & 0.164 & 53 & III \\
13 & N300, G1.5, $\mu$0.85, r=5 & 1.548 & 0.127 & 62 & IV \\
14 & N300, G2.0, $\mu$0.88, r=5 & 1.773 & 0.172 & 54 & II \\
15 & N300, G4.0, $\mu$0.88, r=5 & 1.509 & 0.111 & 63 & I \\
\hline
\end{tabular}
\end{table}

\subsection{Revised Operational Benchmarks}

The original operational threshold of $\eta>0.105$ is formally superseded. Based on
the comprehensive sweep data, three operational tiers are established:

\begin{itemize}
    \item \textbf{Tier 1 (Minimum Viable):} $\eta>1.0$. Achievable at $N\geq150$,
          $G\in[1.0,\,3.0]$, $\mu\geq0.88$.
    \item \textbf{Tier 2 (Deployable):} $\eta>2.0$. Requires $N=300$, $G=1.5$,
          $\mu\geq0.95$.
    \item \textbf{Tier 3 (Optimal):} $\eta>2.5$. Requires $N=300$, $G=1.5$,
          $\mu=0.99$, any $r\in[2,\,10]$.
\end{itemize}

% ── 11. Theoretical Analysis ─────────────────────────────────────────────────
\section{Theoretical Analysis}
\label{sec:theory}

\subsection{Power Law Scaling of $\eta$ with $N$}

The Population Sweep data is consistent with a power law:
\begin{equation}
    \bar{\eta}(N) = \alpha \cdot N^\beta
    \label{eq:powerlaw}
\end{equation}
Fitting to the three data points ($N\in\{50,\,150,\,300\}$) yields exponent
$\beta \approx 0.57$. This sub-linear exponent ($\beta < 1$) is consistent with
the diminishing returns observed in Coordination Quality $Q$: each doubling of $N$
increases $\bar{\eta}$ by less than a factor of 2, reflecting scan-overlap congestion
at the frontier.

The power law enables performance prediction at untested densities. At $N=500$,
extrapolation yields $\bar{\eta}(500) \approx 2.1$ at the Study I parameter settings
($G=4$, $\mu=0.88$). Under the fully-optimized parameters ($G=1.5$, $\mu=0.99$),
the scaling exponent may differ; characterization at $N\in\{500,\,1000\}$ is reserved
for future work.

\subsection{Radius-Invariance Condition}

The scan radius sweep reveals that MAGR is radius-invariant at the optimal
configuration. This can be understood through a coverage density argument. At
$N=300$ on a $50\times50$ grid, the average area per agent is
$2500/300 \approx 8.3$ cells. The Moore neighborhood at $r=2$ covers
$(2\cdot2+1)^2 = 25$ cells. Since $25 > 8.3$, each agent's scan neighborhood
overlaps with the coverage responsibility of multiple neighboring agents at
typical deployment density. The nearest unexplored cell is therefore almost
always within distance $r=2$ of any given agent, making larger scan radii
redundant.

Formally, let $\rho = N / (W \cdot H)$ be the agent surface density. The critical
radius $r^*$ below which radius-invariance breaks down satisfies:
\begin{equation}
    (2r^*+1)^2 \approx \frac{1}{\rho}
    \label{eq:rcrit}
\end{equation}
For $N=300$, $W=H=50$: $r^* \approx \sqrt{1/(4\rho)} - 0.5 \approx 1.4$.
Since $r=2 > r^* \approx 1.4$, all tested radii are above the critical threshold
and therefore equivalent in performance. This analytical prediction is consistent
with the empirical finding that $r=2$ achieves effectively identical efficiency to
$r=10$.

\subsection{Momentum and Trajectory Stability}

The relationship between $\mu$ and $\bar{\eta}$ is monotonically increasing within
the tested range. The KVL update rule:
\begin{equation}
    \vecc{V}_{i}(t+1) = \mu \vecc{V}_{i}(t) + G\cdot\vecc{u}(\theta_i)\cdot 0.12
\end{equation}
can be analyzed as a discrete first-order IIR filter with momentum $\mu$ as the
pole. At $\mu\to 1$, the velocity integrates repulsion impulses over all past
timesteps (infinite memory), producing smooth, directionally stable trajectories.
At $\mu\to 0$, the velocity resets to the current repulsion impulse at each step,
producing maximum responsiveness but erratic, oscillatory paths.

The efficiency improvement from $\mu=0.88$ to $\mu=0.99$ reflects the transition
from a low-memory to a high-memory trajectory regime. In the high-memory regime,
agents maintain sweeping arcs that efficiently cover frontier regions because their
direction is stable over many timesteps. In the low-memory regime, agents change
direction frequently, revisiting cells and producing the redundant search problem
the protocol is designed to eliminate.

% ── 12. Strategic Implementation ─────────────────────────────────────────────
\section{Strategic Implementation Protocol (MVR)}
\label{sec:mvr}

\begin{enumerate}
    \item \textbf{MVR-1: Vector Integration $\geq50$\,Hz.} Sub-threshold frequencies
          introduce Drift-Lag, degrading SSM write accuracy and coverage fidelity.
    \item \textbf{MVR-2: Mesh Latency $<15$\,ms.} Higher latency produces Ghost
          Redundancy --- two agents explore the same cell before the write propagates.
    \item \textbf{MVR-3: Momentum Alignment $\mu_{\text{hw}}\geq0.95$.} The momentum
          sweep establishes that $\mu<0.90$ incurs significant efficiency penalties.
          Hardware implementations must achieve near-unity momentum conservation.
    \item \textbf{MVR-4: Minimum Scan Radius $r=2$.} The radius sweep establishes
          that $r=2$ is sufficient for full MAGR performance. Each agent requires
          only $(2\cdot2+1)^2=25$ SSM reads per timestep.
    \item \textbf{MVR-5: Repulsion Calibration $G\in[1.25,\,1.75]$.} For $N\geq300$
          deployments, repulsion force should be tuned within this range. For
          $N\leq150$ deployments, $G\in[2.50,\,3.00]$ is recommended based on the
          $N\times G$ interaction finding of Study III.
\end{enumerate}

% ── 13. Future Work ───────────────────────────────────────────────────────────
\section{Future Work}
\label{sec:future}

The primary remaining work is physical hardware validation. All results presented
in this paper are derived from simulation on a $50\times50$ occupancy grid.
Deployment on physical robotic platforms will require validation of the five MVR
specifications under real-world conditions, including sensor noise, communication
latency variance, and mechanical constraints not modeled in simulation. In particular,
MVR-3 requires hardware platforms capable of sustaining $\mu\geq0.95$ momentum
conservation in practice, which will constrain the class of viable actuator systems.
Hardware validation constitutes the necessary next step toward operational deployment
of the Vel-v4.3 protocol.

% ── 14. Conclusion ────────────────────────────────────────────────────────────
\section{Conclusion}
\label{sec:conclusion}

This paper presents the most comprehensive empirical study of the Vel decentralized
swarm protocol to date. Five parameter sweeps across 105{,}000 trials characterize
every major protocol dimension and deliver four major findings:

\begin{enumerate}
    \item \textbf{MAGR eliminates boundary pile-up}, restoring monotonic $N$-scaling
          and enabling $N=300$ to outperform $N=150$ by $66.3\%$.
    \item \textbf{The true $G$ optimum is population-dependent}: $G=1.5$ at $N=300$
          and $G=2.75$ at $N=150$, due to a $N\times G$ interaction effect in which
          higher density requires lower repulsion to avoid cell overshoot.
    \item \textbf{The original $\mu=0.88$ assumption was suboptimal by $49.6\%$.}
          $\mu=0.99$ simultaneously maximizes efficiency ($\bar{\eta}=2.696$) and
          minimizes variance ($\sigma=0.089$), occupying the dominant position on the
          efficiency--stability frontier.
    \item \textbf{MAGR is radius-invariant.} $\Delta\eta=0.006$ across $r\in[2,\,10]$,
          while computational cost scales $17.6\times$. The optimal deployment radius
          is $r=2$, analytically explained by the critical radius condition
          $(2r^*+1)^2 \approx 1/\rho$.
\end{enumerate}

The global optimum --- $N=300$, $G=1.5$, $\mu=0.99$, $r=2$, $\bar{\eta}=2.696$ ---
represents a $78.7\%$ efficiency gain over the initial Study I baseline. The revised
operational framework establishes three deployment tiers ($\eta>1.0$, $\eta>2.0$,
$\eta>2.5$) superseding the original single threshold of $\eta>0.105$ by a factor
of $19\times$.

% ── Acknowledgements ──────────────────────────────────────────────────────────
\section*{Acknowledgements}
The author acknowledges Gnaneswar Kumar Pullaganti and Dakshith Karthikeyan for their
contributions during the early science fair phase of the Vel project. Their involvement
was limited to that initial stage; they had no role in the development of Vel-v4.3,
the MAGR architecture, the parameter sweeps reported in this paper, or any part of
this work.

The author used Claude (Anthropic) as an AI writing assistant during the preparation
of this manuscript. All research concepts, experimental design, simulation implementation,
data collection, and scientific conclusions are the original work of the author. The AI
assistant was used for LaTeX formatting and prose editing only.

The author's ORCID is 0009-0006-2907-8815.

% ── References ────────────────────────────────────────────────────────────────
\begin{thebibliography}{1}

\bibitem{dorigo2004}
M.~Dorigo and E.~\c{S}ahin, ``Swarm robotics: Special issue,''
\emph{Autonomous Robots}, vol.~17, pp.~111--113, 2004.

\bibitem{brambilla2013}
M.~Brambilla, E.~Ferrante, M.~Birattari, and M.~Dorigo,
``Swarm robotics: A review from the swarm engineering perspective,''
\emph{Swarm Intelligence}, vol.~7, no.~1, pp.~1--41, 2013.

\bibitem{jamshidpey2024}
A.~Jamshidpey et al., ``Centralization vs.\ decentralization in multi-robot
sweep coverage with ground robots and {UAV}s,''
\emph{arXiv preprint arXiv:2408.06553}, 2024.

\bibitem{moravec1985}
H.~P.~Moravec and A.~Elfes, ``High resolution maps from wide angle sonar,''
in \emph{Proc. IEEE ICRA}, 1985, pp.~116--121.

\bibitem{elfes1989}
A.~Elfes, ``Using occupancy grids for mobile robot perception and navigation,''
\emph{Computer}, vol.~22, no.~6, pp.~46--57, 1989.

\bibitem{reynolds1987}
C.~W.~Reynolds, ``Flocks, herds and schools: A distributed behavioral model,''
in \emph{Proc. SIGGRAPH '87}, 1987, pp.~25--34.

\bibitem{kennedy1995}
J.~Kennedy and R.~C.~Eberhart, ``Particle swarm optimization,''
in \emph{Proc. IEEE Int. Conf. Neural Networks}, 1995, pp.~1942--1948.

\bibitem{dorigo1992}
M.~Dorigo, V.~Maniezzo, and A.~Colorni, ``Ant system: Optimization by a colony
of cooperating agents,''
\emph{IEEE Trans. Syst., Man, Cybern. B}, vol.~26, no.~1, pp.~29--41, 1996.

\bibitem{galceran2013}
E.~Galceran and M.~Carreras, ``A survey on coverage path planning for robotics,''
\emph{Robotics and Autonomous Systems}, vol.~61, no.~12, pp.~1258--1276, 2013.

\bibitem{li2024}
Q.~Li et al., ``Improved ant colony optimization algorithm based on islands type
for mobile robot path planning,''
\emph{International Journal of Advanced Robotic Systems}, 2024.

\bibitem{harris2020}
C.~R.~Harris et al., ``Array programming with {NumPy},''
\emph{Nature}, vol.~585, pp.~357--362, 2020.

\bibitem{hunter2007}
J.~D.~Hunter, ``Matplotlib: A 2{D} graphics environment,''
\emph{Computing in Science \& Engineering}, vol.~9, no.~3, pp.~90--95, 2007.

\end{thebibliography}

% --- Author Biography Section ---
\section*{Author Biography}
% The [-2ex] reduces the gap between the last section and the biography
\vspace{-7ex} 

\begin{IEEEbiographynophoto}{Dakshith Rajkumar}
is an independent researcher specializing in decentralized swarm intelligence and autonomous robotic navigation. His work focuses on the development of the Vel protocol series, specifically exploring Map-Aware Gradient Repulsion (MAGR) and spatio-temporal efficiency in multi-agent systems. A member of the National Junior Honor Society (NJHS) at Fowler Middle School, part of the Frisco Independent School District, he maintains a rigorous commitment to academic excellence and emergent engineering logic.
\end{IEEEbiographynophoto}

\vfill % Pushes the bio to the bottom or keeps spacing consistent with IEEE style

\end{document}